% Generated by `taskgen course examples` from powmod.montgomery (seed 4).
% Do not edit: change examples.yaml and regenerate.
\begin{example}[возведение в степень методом Монтгомери]\label{ex:exp-montgomery}
Возвести целое число в степень по модулю: \[13^{98}\ (\textrm{mod}\ 331)\] с помощью метода Монтгомери. Обосновать выбор параметра $R$.

\par\noindent\textit{Решение.} $R = 2^{9} = 512$ "--- наименьшая степень двойки, превосходящая модуль $331$; $R$ взаимно прост с нечётным модулем, а деление на $R$ "--- сдвиг. $m' = -m^{-1} \bmod R = 413$, $R^2 \bmod m = 323$ (здесь $m = 331$).
 Переводим $a$ в форму Монтгомери $\bar a = \mathrm{REDC}(a \cdot (R^2 \bmod m)) = aR \bmod m$; далее для битов показателя $1100010_2$ после старшего: $\bar x \gets \mathrm{REDC}(\bar x^2)$, при бите $1$ ещё $\bar x \gets \mathrm{REDC}(\bar x \bar a)$; в конце $x = \mathrm{REDC}(\bar x)$.\par $\mathrm{REDC}(T)$: $u = (T \bmod R)\, m' \bmod R$, $t = (T + u m)/R$, при $t \ge m$ $t \gets t - m$.
\[\begin{array}{l|r|r|r|r}
 & T & u & t & \mathrm{REDC}(T) \\ \hline
\bar a = a \cdot R^2 & 4199 & 43 & 36 & 36 \\
\text{бит 1}: \bar x^2 & 1296 & 208 & 137 & 137 \\
\bar x \cdot \bar a & 4932 & 180 & 126 & 126 \\
\text{бит 0}: \bar x^2 & 15876 & 116 & 106 & 106 \\
\text{бит 0}: \bar x^2 & 11236 & 212 & 159 & 159 \\
\text{бит 0}: \bar x^2 & 25281 & 349 & 275 & 275 \\
\text{бит 1}: \bar x^2 & 75625 & 101 & 213 & 213 \\
\bar x \cdot \bar a & 7668 & 164 & 121 & 121 \\
\text{бит 0}: \bar x^2 & 14641 & 13 & 37 & 37 \\
\bar x \to x & 37 & 433 & 280 & 280 \\
\end{array}\]
\[13^{98} \equiv 280\ (\textrm{mod}\ 331).\]
\end{example}
